\documentclass[11pt,a4paper]{article}

\usepackage[a4paper,margin=2.5cm]{geometry}
\usepackage{fontspec}
\usepackage{titlesec}
\usepackage{tikz}
\usetikzlibrary{positioning,arrows.meta,fit,backgrounds}
\usepackage{xcolor}
\usepackage{hyperref}
\usepackage{enumitem}
\usepackage{fancyhdr}
\usepackage{tcolorbox}
\tcbuselibrary{skins,breakable}

\setmainfont{Sylfaen}
\newfontfamily\latinfont{Georgia}

\definecolor{brandgreen}{HTML}{1F6F5C}
\definecolor{brandlight}{HTML}{E7F1EE}
\definecolor{boxframe}{HTML}{1F6F5C}

\hypersetup{
  colorlinks=true,
  linkcolor=brandgreen,
  urlcolor=brandgreen,
  pdftitle={LLM-Climate-Health -- ტექნიკური დოკუმენტაცია},
  pdfauthor={Sergi Koniashvili}
}

\titleformat{\section}{\large\bfseries\color{brandgreen}}{\thesection}{0.5em}{}
\titleformat{\subsection}{\normalsize\bfseries}{\thesubsection}{0.5em}{}

\setlength{\headheight}{14pt}
\pagestyle{fancy}
\fancyhf{}
\fancyhead[L]{LLM-Climate-Health}
\fancyhead[R]{ტექნიკური დოკუმენტაცია}
\fancyfoot[C]{\thepage}

\newtcolorbox{infobox}[1]{
  colback=brandlight, colframe=boxframe, boxrule=0.6pt, arc=2pt,
  breakable, title={#1}, fonttitle=\bfseries, coltitle=white,
  colbacktitle=brandgreen
}

\begin{document}

\begin{titlepage}
\centering
\vspace*{3cm}
{\Huge\bfseries\color{brandgreen} LLM-Climate-Health}\\[0.5cm]
{\Large ტექნიკური დოკუმენტაცია და ლექსიკონი}\\[2cm]
{\large LLM-ზე დაფუძნებული ETL პლატფორმა კლიმატისა და\\ ჯანმრთელობის მონაცემთა ინტეგრაციისთვის}\\[2cm]
{\large საბაკალავრო ნაშრომი}\\[0.3cm]
{\large ბრემენის უნივერსიტეტი -- ინფორმატიკის ფაკულტეტი}\\[2cm]
{\large ავტორი: \textbf{სერგი კონიაშვილი}}\\[0.5cm]
{\large თარიღი: \today}\\[3cm]
{\normalsize ეს დოკუმენტი აღწერს პლატფორმის არქიტექტურას, მონაცემთა\\
ინტეგრაციის პროცესს და მასში გამოყენებულ ძირითად ცნებებს --\\
საბაკალავრო ნაშრომში გამოსაყენებლად.}
\vfill
\end{titlepage}

\tableofcontents
\newpage

%=====================================================================
\section{შესავალი}

ბევრი დაავადება კლიმატზეა დამოკიდებული -- მათი გადამტანები (მაგალითად,
კოღოები) მგრძნობიარენი არიან ტემპერატურის, ნალექისა და ტენიანობის მიმართ.
ეს ნიშნავს, რომ კლიმატური მონაცემები შეიძლება გამოყენებულ იქნას როგორც
ადრეული ინდიკატორი დაავადების აფეთქების რისკისთვის.

ამ პროექტის მიზანია შეიქმნას ზოგადი დანიშნულების პლატფორმა, რომელიც
საშუალებას მისცემს მომხმარებელს -- პროგრამირების ან კლიმატოლოგიის
ცოდნის გარეშე -- ბუნებრივი ენით ჩამოაყალიბოს მოთხოვნა (დაავადება,
რეგიონი, კლიმატური ცვლადები) და მიიღოს ანალიზისთვის მზა, ურთიერთდაკავშირებული
მონაცემები. დიდი ენობრივი მოდელი (LLM) ამ მოთხოვნას გარდაქმნის
შესრულებად ინტეგრაციის გეგმად და ყოველ ნაბიჯს მარტივი ენით უხსნის
მომხმარებელს.

\textbf{მნიშვნელოვანია}: პლატფორმა აწარმოებს \emph{მონაცემებს}, არა
\emph{პროგნოზებს}. შედეგი გამოიყენება როგორც შემავალი მონაცემები
შემდგომი დაავადების მოდელებისთვის.

\begin{infobox}{ვალიდაციის სცენარი}
პროექტი ვალიდირებულია დენგეს ცხელების მაგალითზე ტაილანდში, თუმცა
არქიტექტურა შექმნილია ისე, რომ იგივე ლოგიკა ვრცელდებოდეს სხვა
ვექტორგადამტან დაავადებებზეც (მალარია, ქოლერა, ლაიმის დაავადება) და
ნებისმიერ რეგიონზე, სადაც არსებობს საკმარისი მონაცემები.
\end{infobox}

%=====================================================================
\section{სისტემის არქიტექტურა}

პლატფორმა აგებულია სამი ფენისგან: \textbf{frontend} (Flutter),
\textbf{backend} (Python/FastAPI) და \textbf{გარე მონაცემთა წყაროები}.
Flutter არჩეულია იმიტომ, რომ ერთი კოდით შესაძლებელია აპლიკაციის აგება
ვებისთვის, მობილურისთვის და დესკტოპისთვისაც -- ეს მნიშვნელოვანია, რადგან
პლატფორმა გამიზნულია ფართო ხელმისაწვდომობისთვის ბმულის საშუალებით.

\begin{figure}[h!]
\centering
\begin{tikzpicture}[
  box/.style={draw=boxframe, fill=brandlight, rounded corners=3pt, minimum width=4.2cm, minimum height=1cm, align=center, font=\small},
  arr/.style={-{Stealth[length=2.5mm]}, thick, color=boxframe}
]
\node[box, fill=brandgreen, text=white] (fe) {Flutter Frontend\\ (Web / Mobile / Desktop)};
\node[box, below=1.3cm of fe] (be) {FastAPI Backend\\ (ვალიდაცია, ETL, cache, rate-limit)};

\node[box, below=1.6cm of be, xshift=-4.6cm, minimum width=2.8cm] (clim) {კლიმატის API-ები\\ Open-Meteo / NASA POWER};
\node[box, below=1.6cm of be, minimum width=2.8cm] (dis) {დაავადების\\ არქივი OpenDengue};
\node[box, below=1.6cm of be, xshift=4.6cm, minimum width=2.8cm] (custom) {მომხმარებლის\\ URL / ატვირთვა};
\node[box, below=3.6cm of be, minimum width=2.8cm] (llm) {Gemini LLM\\ (ახსნა ჩვეულებრივ ენაზე)};

\draw[arr] (fe) -- (be) node[midway, right, font=\scriptsize] {REST / JSON};
\draw[arr] (be) -- (clim);
\draw[arr] (be) -- (dis);
\draw[arr] (be) -- (custom);
\draw[arr] (be) -- (llm);
\end{tikzpicture}
\caption{პლატფორმის მაღალი დონის არქიტექტურა}
\end{figure}

\subsection{Backend-ის მოდულები}
\begin{itemize}[leftmargin=1.4em]
  \item \textbf{config} -- გარემოს ცვლადები, მხარდაჭერილი რეგიონების/წყაროების რეესტრი;
  \item \textbf{services/dengue} -- დაავადების შემთხვევების მონაცემები (OpenDengue, URL, ატვირთვა);
  \item \textbf{services/climate} -- კლიმატური რეანალიზის მონაცემები (Open-Meteo, NASA POWER);
  \item \textbf{services/etl} -- ორკესტრაცია: მოძიება, აგრეგირება, დაკავშირება;
  \item \textbf{services/llm} -- Gemini API-სთან კომუნიკაცია, rate-limiting;
  \item \textbf{services/cache} -- შედეგების დროებითი შენახვა და \emph{last\_verified} დროის ნიშნულის დამატება.
\end{itemize}

%=====================================================================
\section{მონაცემთა ინტეგრაციის პროცესი (ETL)}

ყოველი მოთხოვნა გადის შემდეგ ეტაპებს, რომლებიც მომხმარებელს ჩანს
ეკრანზე -- ეს იმისთვისაა, რომ მომხმარებელმა დარწმუნებით იცოდეს, რომ
მისი ბრძანება სწორად შესრულდა.

\begin{figure}[h!]
\centering
\begin{tikzpicture}[
  etlstep/.style={draw=boxframe, fill=white, rounded corners=3pt, text width=9.5cm, minimum height=0.9cm, align=left, font=\small, inner sep=6pt},
  arr/.style={-{Stealth[length=2.5mm]}, thick, color=boxframe}
]
\node[etlstep] (s1) {\textbf{1.} მოთხოვნის მიღება და ვალიდაცია};
\node[etlstep, below=0.35cm of s1] (s2) {\textbf{2.} რეგიონის მიბმა საცნობარო კოორდინატებზე};
\node[etlstep, below=0.35cm of s2] (s3) {\textbf{3.} დაავადების შემთხვევების მონაცემების მოძიება};
\node[etlstep, below=0.35cm of s3] (s4) {\textbf{4.} კლიმატური მონაცემების მოძიება};
\node[etlstep, below=0.35cm of s4] (s5) {\textbf{5.} დღიური მონაცემების თვიურ სიხშირემდე აგრეგირება};
\node[etlstep, below=0.35cm of s5] (s6) {\textbf{6.} ორივე სერიის დროში დაკავშირება (join)};
\node[etlstep, below=0.35cm of s6] (s7) {\textbf{7.} ახსნის გენერირება (LLM) და პასუხის დაბრუნება};

\draw[arr] (s1) -- (s2);
\draw[arr] (s2) -- (s3);
\draw[arr] (s3) -- (s4);
\draw[arr] (s4) -- (s5);
\draw[arr] (s5) -- (s6);
\draw[arr] (s6) -- (s7);
\end{tikzpicture}
\caption{ETL პროცესის ეტაპები (ერთი მოთხოვნისთვის)}
\end{figure}

\subsection{გამჭვირვალობა და კეშირება}
თითოეული ნაბიჯი ტექსტურად ფორმდება და უბრუნდება frontend-ს, ისე რომ
მომხმარებელს შეუძლია ნახოს ზუსტად რა მოხდა. თუ იგივე მოთხოვნა მეორდება
15 წუთის განმავლობაში, სისტემა უბრუნებს კეშირებულ შედეგს (გარე API-ების
ზედმეტი დატვირთვის თავიდან ასაცილებლად) და უჩვენებს \emph{ბოლო
გადამოწმების} დროის ნიშნულს (\texttt{last\_verified}), რომ მომხმარებელმა
იცოდეს რამდენად ახალია მონაცემი.

ახსნის ტექსტი (\texttt{explanation}) ყოველთვის Gemini-ს მეშვეობით
გენერირდება, მაგრამ თუ API-ის გამოძახება ჩავარდება (მაგ. rate limit ან
დროებითი მიუწვდომლობა), სისტემა წინასწარ მომზადებულ, ზოგად ტექსტს
აბრუნებს ისე, რომ მომხმარებელმა საერთოდ არაფერი არ დაინახოს ცარიელი
ადგილის ნაცვლად. რომ ეს ჩანაცვლება არასდროს გავიდეს ნამდვილ,
კონკრეტულ მოთხოვნაზე მორგებულ ტექსტად, პასუხს დაემატა ცალკე ველი
\texttt{explanation\_source} (\texttt{"llm"} ან \texttt{"fallback"}) --
frontend პირდაპირ აჩვენებს, რომელი შემთხვევა მოხდა.

%=====================================================================
\section{ბუნებრივი ენით მოთხოვნის ინტერპრეტაცია}
\label{sec:nl-parsing}

თავდაპირველად მომხმარებელი მოთხოვნას მხოლოდ სტრუქტურირებული ფორმის
(dropdown-ების და თარიღის ველების) მეშვეობით ავსებდა. ამას დაემატა
დამატებითი, არასავალდებულო გზა: მომხმარებელს შეუძლია მოთხოვნა
თავისუფალი ტექსტით აღწეროს (მაგ. ,,დენგე ტაილანდში 2015-დან 2023
წლამდე, ყოველთვიურად, ტემპერატურა და ნალექი''), და Gemini LLM ცდილობს
ამის საფუძველზე ფორმის ველების ავტომატურად შევსებას
(\texttt{POST /api/parse-request}).

\begin{infobox}{უსაფრთხოების პრინციპი: მხოლოდ შევსება, არასდროს გაშვება}
ბუნებრივი ენის პარსინგი მხოლოდ წინასწარ ავსებს ფორმას -- ის არასდროს
უშვებს ETL პროცესს პირდაპირ. მომხმარებელს ყოველთვის შეუძლია შევსებული
ველების გადამოწმება და საჭიროების შემთხვევაში ხელით შესწორება, სანამ
,,Run integration''-ზე დააჭერს. ეს ინარჩუნებს პლატფორმის ძირითად
პრინციპს: LLM ხსნის და ეხმარება, მაგრამ არასდროს იღებს გადაწყვეტილებას
მომხმარებლის ნაცვლად.
\end{infobox}

\subsection{დახურული ლექსიკონი გამონაგონის თავიდან ასაცილებლად}
LLM-ს არ აქვს თავისუფლება ნებისმიერი მნიშვნელობის შემოთავაზების --
prompt-ში პირდაპირ ეწერება ყველა დაშვებული დაავადება, რეგიონი, ცვლადი,
აგრეგაცია და კლიმატის წყარო, და მოდელს მოეთხოვება აირჩიოს მხოლოდ ამ
სიებიდან, ან დატოვოს ველი ცარიელი (\texttt{null}), თუ ტექსტში ცალსახად
არ ჩანს. პასუხის დამუშავებისას (\texttt{\_coerce\_parsed\_json})
ყოველი მნიშვნელობა ხელახლა მოწმდება backend-ის კოდში -- თუ მოდელმა
მაინც შემოგვთავაზა ამ სიების გარეთა მნიშვნელობა, ის იშლება
(\texttt{None}-ად იქცევა) და ეს ფაქტი მომხმარებელს უჩნდება
,,notes'' ველში ახსნილი დაშვების სახით. ამგვარად, არასწორად
ამოცნობილი დაავადება ან რეგიონი ვერასდროს მიაღწევს პირდაპირ ETL
პროცესამდე.

%=====================================================================
\section{საკვანძო ცნებები კლიმატურ მონაცემებში}

ეს თავი გამოსადეგია, როგორც მოკლე სასწავლო მასალა -- განმარტავს
ტერმინებს, რომლებსაც პლატფორმა იყენებს.

\subsection{რეანალიზი (Reanalysis)}
რეანალიზი არის წარსული ამინდის მონაცემების მეცნიერულად აღდგენილი,
თანმიმდევრული სურათი. ის იქმნება ისტორიული დაკვირვებების (მეტეოსადგურები,
თანამგზავრები, ხომალდები) და ფიზიკური ამინდის მოდელის შერწყმით. ეს
\emph{არ არის} ერთი კონკრეტული სადგურის ნედლი მონაცემი -- ის წარმოადგენს
საუკეთესო შეფასებას იმისა, თუ როგორი იყო ატმოსფერო კონკრეტულ ადგილას და
დროს. პლატფორმა იყენებს ECMWF ERA5 რეანალიზს (Open-Meteo-ს საშუალებით)
და NASA POWER-ს (MERRA-2 რეანალიზზე დაფუძნებული).

\subsection{რეანალიზი vs. მოდელის პროგნოზი (Projection)}
რეანალიზი აღწერს \emph{წარსულს} და კორექტირდება რეალურ დაკვირვებებზე.
მოდელის პროგნოზი (მაგ. CMIP6/CORDEX სცენარები) კი სიმულირებს
\emph{მომავალ} კლიმატს ემისიების სხვადასხვა დაშვებით და არ არის
დაკავშირებული კონკრეტულ დაკვირვებულ მოვლენებთან. ამჟამინდელი
იმპლემენტაცია იყენებს მხოლოდ ისტორიულ რეანალიზს (\texttt{scenario=historical});
მომავალი პროგნოზების ინტეგრაცია დაგეგმილია შემდეგ ეტაპზე.

\subsection{Bias Correction (მიკერძოების კორექცია)}
კლიმატის მოდელებს ხშირად აქვთ სისტემატური გადახრა რეალობისგან
(მაგალითად, მუდმივად აჭარბებენ ნალექის რაოდენობას კონკრეტულ რეგიონში).
Bias correction არის სტატისტიკური მეთოდი, რომელიც ასწორებს ამ გადახრას
ისტორიულ დაკვირვებებთან შედარებით, სანამ მონაცემი გამოყენებული იქნება
შემდგომი ანალიზისთვის. ეს განსაკუთრებით მნიშვნელოვანია მაშინ, როცა
პროგნოზული (არა-რეანალიზის) მონაცემები გამოიყენება.

\subsection{ემისიის სცენარები (Emission Scenarios)}
ეს არის დაშვებები მომავალი სათბურის გაზების გამონაბოლქვის შესახებ
(მაგ. SSP1-2.6 -- დაბალი ემისია, SSP5-8.5 -- მაღალი ემისია), რომლებზეც
დაფუძნებულია კლიმატის მომავალი პროგნოზები. რადგან მომავალი ზუსტად
უცნობია, მეცნიერები იყენებენ რამდენიმე სცენარს პარალელურად.

\subsection{სივრცული და დროითი შეთანხმება (Join)}
სხვადასხვა წყაროს მონაცემები განსხვავებული სივრცული (წერტილი, ბადე,
რეგიონი) და დროითი (დღიური, თვიური, წლიური) გარჩევადობით მოდის.
პლატფორმა კლიმატურ მონაცემებს დღიურიდან თვიურამდე აგრეგირებს
(ტემპერატურა -- საშუალო, ნალექი -- ჯამი), რათა ის დაემთხვეს დაავადების
შემთხვევების ჩვეულებრივ თვიურ ანგარიშგებას, შემდეგ კი აერთიანებს
ორივეს საერთო თვე-წლის გასაღებით.

%=====================================================================
\section{გამოყენებული მონაცემთა წყაროები}

\begin{itemize}[leftmargin=1.4em]
  \item \textbf{OpenDengue} -- აკადემიური, საჯაროდ ხელმისაწვდომი დენგეს
    შემთხვევების არქივი (Clarke et al., 2024, \emph{Scientific Data}).
    მოიცავს 39 ქვეყანას საკმარისი (12+ თვე) თვიური გარჩევადობით.
  \item \textbf{Open-Meteo Historical Archive} -- ECMWF ERA5 რეანალიზზე
    დაფუძნებული, უფასო, გასაღების გარეშე ხელმისაწვდომი API.
  \item \textbf{NASA POWER} -- ალტერნატიული სამეცნიერო კლიმატური წყარო
    (MERRA-2 რეანალიზი), ასევე გასაღების გარეშე ხელმისაწვდომი.
  \item \textbf{Thai Meteorological Department (TMD)} -- ტაილანდის
    ეროვნული, ოფიციალური მეტეოროლოგიური სამსახურის API
    (\texttt{data.tmd.go.th}), ხელმისაწვდომი მხოლოდ ტაილანდისთვის
    (იხ. ქვემოთ).
  \item \textbf{მომხმარებლის საკუთარი წყარო} -- URL ან ატვირთული CSV
    ფაილი, თუ მომხმარებელს სურს საკუთარი სამეცნიერო მონაცემების გამოყენება.
\end{itemize}

\subsection{რეგიონულად შეზღუდული წყარო: TMD}
ღია, გლობალური კლიმატის წყაროებისგან (Open-Meteo, NASA POWER)
განსხვავებით, TMD მხოლოდ ტაილანდის ტერიტორიაზეა ხელმისაწვდომი და
მოითხოვს რეგისტრაციას (\texttt{uid}/\texttt{ukey} წყვილი, უფასო,
\texttt{data.tmd.go.th}-ზე). ეს შეზღუდვა კოდში ცალკე რეესტრშია
აღწერილი (\texttt{CLIMATE\_SOURCE\_REGIONS}), რომელსაც backend
ამოწმებს მოთხოვნის დროს -- თუ მომხმარებელი TMD-ს სხვა რეგიონისთვის
აირჩევს, მკაფიო 400 შეცდომას იღებს დაბნეული ან ცარიელი პასუხის
ნაცვლად.

\begin{infobox}{გულწრფელი შეზღუდვა: დაუდასტურებელი პასუხის სტრუქტურა}
TMD-ს საჯარო დოკუმენტაცია აღწერს endpoint-ებს და საჭირო
\texttt{uid}/\texttt{ukey} პარამეტრებს, მაგრამ არ აქვეყნებს ზუსტ JSON
ველების სახელებს. შესაბამისად, პასუხის დამუშავების ლოგიკა
(\texttt{\_parse\_tmd\_response}) დაწერილია დოკუმენტაციაში აღწერილი
ველების საფუძველზე და საჭიროებს რეალურ smoke-ტესტს, როგორც კი
წვდომის მონაცემები (\texttt{TMD\_API\_UID}/\texttt{TMD\_API\_UKEY})
დარეგისტრირდება. ეს გაურკვევლობა შეგნებულად არის იზოლირებული ერთ
ფუნქციაში, რომ საჭიროების შემთხვევაში გასწორება ერთ ადგილას
შემოიფარგლოს.
\end{infobox}

%=====================================================================
\section{დაავადების მონაცემთა წყაროების ცოცხალი ძებნა}
\label{sec:source-discovery}

თავდაპირველ ვერსიაში თითო დაავადებას მხოლოდ \emph{ერთი}, კოდში
ჩაშენებული წყარო ჰქონდა მიბმული -- მომხმარებელს არჩევანის საშუალება არ
ჰქონდა. ამის გამოსასწორებლად დაემატა ცოცხალი ძებნის ფუნქცია: backend-ი
რეალურ დროში ეკითხება ორ სანდო, საჯარო, გასაღების გარეშე ხელმისაწვდომ
სამეცნიერო/ოფიციალურ კატალოგს და მომხმარებელს რეალურ, დამოწმებად
შედეგებს სთავაზობს ასარჩევად.

\begin{infobox}{პრინციპი: არაფერი გამოგონილი}
დიდი ენობრივი მოდელი (LLM) აქ \emph{არ} გამოიყენება წყაროს მოსაძებნად ან
URL-ის ან citation-ის შესათხზად -- ეს საფრთხეს შეუქმნიდა მონაცემის
სანდოობას. ძებნა ხდება მხოლოდ ორი კონკრეტული, სანდო API-ის მეშვეობით;
დაბრუნებული ყოველი ველი (სათაური, ორგანიზაცია, ბმული) პირდაპირ
წამოღებულია იმ API-ის საკუთარი პასუხიდან.
\end{infobox}

\subsection{ორი წყარო-კატალოგი}
\begin{itemize}[leftmargin=1.4em]
  \item \textbf{WHO Global Health Observatory (GHO) API} -- ჯანმო-ს
    საკუთარი, ოფიციალური OData API (\texttt{ghoapi.azureedge.net}).
    დაავადების საკვანძო სიტყვით (მაგ. ,,Malaria'') ძებნა აბრუნებს
    ათეულობით რეალურ ინდიკატორს (,,დადასტურებული შემთხვევები'',
    ,,შეფასებული შემთხვევები'', ,,სიკვდილიანობა'' და ა.შ.) -- მომხმარებელი
    თავად ირჩევს რომელი ზუსტად ჰინტერესებს, ერთადერთი ჩაშენებული
    ინდიკატორის ნაცვლად.
  \item \textbf{HDX -- Humanitarian Data Exchange} (\texttt{data.humdata.org})
    -- გაერო-ს (UN OCHA) კატალოგი, სადაც სამთავრობო/ჯანმო-ს/არასამთავრობო
    ორგანიზაციები პირდაპირ ატვირთავენ ჩამოსატვირთ CSV/XLSX ფაილებს.
    გამოსადეგია მაშინ, როცა WHO GHO-ს არ აქვს საკმარისად დეტალური
    (მაგ. თვიური) მონაცემი კონკრეტული დაავადებისთვის.
\end{itemize}

\subsection{როგორ ჯდება არსებულ არქიტექტურაში}
მთავარი დიზაინის გადაწყვეტილება: HDX-დან ნაპოვნი შედეგი \emph{არ}
საჭიროებს ახალ parser-ს -- ის უბრალოდ გადის უკვე არსებული და
დატესტილი \texttt{custom\_url} channel-ით (\texttt{get\_case\_data\_from\_url}),
რადგან HDX პირდაპირ ჩამოსატვირთ ცხრილოვან ფაილებს გვთავაზობს. მხოლოდ
WHO GHO-სთვის დაემატა ახალი, სპეციალური ადაპტერი
(\texttt{services/who\_gho.py}), რადგან მისი მონაცემები ცხრილის ნაცვლად
სტრუქტურირებული API-დანაა წამოსაღები.

WHO GHO-ს ინდიკატორები ყოველთვის წლიურია, მაშინაც კი, როცა თავად
დაავადებას (მაგ. დენგე OpenDengue-ს მეშვეობით) თვიური მონაცემი აქვს.
ETL პროცესი (\texttt{run\_integration}) ამ შემთხვევაში მიმართვის
გარჩევადობას (\texttt{resolution}) დაავადების ნაგულისხმევი მნიშვნელობის
ნაცვლად რეალურად გამოძახებული წყაროდან იღებს -- წინააღმდეგ შემთხვევაში
დაავადების და კლიმატის სერიები არასწორ, შეუთავსებელ პერიოდებზე
შეიძლებოდა გაერთიანებულიყო.

\begin{figure}[h!]
\centering
\begin{tikzpicture}[
  box/.style={draw=boxframe, fill=brandlight, rounded corners=3pt, minimum width=4.4cm, minimum height=1cm, align=center, font=\small},
  arr/.style={-{Stealth[length=2.5mm]}, thick, color=boxframe}
]
\node[box, fill=brandgreen, text=white] (req) {მომხმარებელი: დაავადება + რეგიონი};
\node[box, below=1.3cm of req] (search) {\texttt{GET /api/search-case-sources}};

\node[box, below=1.5cm of search, xshift=-2.6cm] (who) {WHO GHO API\\ (ინდიკატორების ძებნა)};
\node[box, below=1.5cm of search, xshift=2.6cm] (hdx) {HDX API\\ (dataset-ების ძებნა)};

\node[box, below=1.5cm of who, xshift=2.6cm, fill=brandgreen, text=white] (pick) {მომხმარებელი ირჩევს ერთს};

\node[box, below=1.3cm of pick, xshift=-2.6cm] (whoingest) {ახალი ადაპტერი:\\ \texttt{who\_gho.py}};
\node[box, below=1.3cm of pick, xshift=2.6cm] (urlingest) {არსებული channel:\\ \texttt{custom\_url}};

\draw[arr] (req) -- (search);
\draw[arr] (search) -- (who);
\draw[arr] (search) -- (hdx);
\draw[arr] (who) -- (pick);
\draw[arr] (hdx) -- (pick);
\draw[arr] (pick) -- (whoingest);
\draw[arr] (pick) -- (urlingest);
\end{tikzpicture}
\caption{წყაროების ცოცხალი ძებნის ნაკადი}
\end{figure}

არჩეული წყაროს citation ავტომატურად ჩნდება საბოლოო შედეგის ,,Sources''
სექციაში, ისე რომ მომხმარებელს ყოველთვის ზუსტად ეცოდინება საიდან
მოვიდა თითოეული რიცხვი.

\subsection{შეზღუდვები}
\begin{itemize}[leftmargin=1.4em]
  \item ძებნა შემოიფარგლება მხოლოდ ამ ორი კატალოგით -- ეს შეგნებული
    არჩევანია (ვიწრო, სანდო allowlist ღია web-ძებნის ნაცვლად), მაგრამ
    ნიშნავს, რომ ამ ორ წყაროში არარსებული მონაცემი უბრალოდ არ გამოჩნდება;
  \item WHO GHO-ს ინდიკატორები თითქმის ყოველთვის წლიურია -- უფრო
    დეტალური (თვიური) ანალიზისთვის HDX ან საკუთარი URL/ფაილი რჩება
    სჯობს არჩევანად;
  \item ორივე API საჯარო და უფასოა, მაგრამ გარე სერვისებია -- დროებითი
    მიუწვდომლობის შემთხვევაში ძებნა უბრალოდ ცარიელ სიას აბრუნებს
    (არა შეცდომას), და მომხმარებელს კვლავ შეუძლია custom URL/ატვირთვა.
\end{itemize}

%=====================================================================
\section{მონაცემთა ინტელექტუალური ჰარმონიზაცია და გამჭვირვალობის ჟურნალი}
\label{sec:data-harmonization}

სამეცნიერო და ჰუმანიტარულ მონაცემთა ბაზებში (მაგ. HDX, კვლევითი ცენტრების
CSV ფაილები) სვეტების სახელები და ფორმატები ხშირად არასტანდარტულია
(მაგ. \texttt{period}, \texttt{observation\_date}, \texttt{malaria\_cases},
\texttt{incidence\_rate}). 

პლატფორმაში დაინერგა \textbf{AI Data Harmonizer \& Audit System}:
\begin{itemize}[leftmargin=1.4em]
  \item \textbf{სვეტების მოქნილი ამოცნობა და Fuzzy Matching}: სისტემა
    ავტომატურად აიდენტიფიცირებს დროითი ინდექსის სვეტს და დაავადების
    შემთხვევების მეტრიკას, მაშინაც კი თუ მათი დასახელებები განსხვავდება;
  \item \textbf{მონაცემთა გაწმენდა და ფორმატირება}: არასწორი ან ცარიელი
    ჩანაწერების გაფილტვრა რეალური მონაცემების დაუმახინჯებლად;
  \item \textbf{გამჭვირვალობის ჟურნალი (Data Transformation Audit Log)}:
    პროგრამა აღრიცხავს ყველა ცვლილებას -- რა სვეტები მიიღო, რა გარდაქმნა,
    რამდენი სტრიქონი გაიფილტრა და რა დარჩა;
  \item \textbf{AI-ს მიერ გენერირებული ადამიანური ახსნა}: Gemini LLM
    მომხმარებელს მარტივ ენაზე უხსნის მონაცემთა დამუშავების მთლიან პროცესს,
    რაც თავიდან არიდებს გაუგებრობებსა და მონაცემთა არასწორ ინტერპრეტაციას.
\end{itemize}

%=====================================================================
\section{მრავალგანზომილებიანი ინტერაქტიული ვიზუალიზაცია და KPI მეტრიკები}
\label{sec:interactive-visualization}

კლიმატისა და ეპიდემიოლოგიური მონაცემების სიღრმისეული ანალიზისთვის
ინტერფეისს დაემატა მართვადი ვიზუალიზაციის კომპონენტები:
\begin{itemize}[leftmargin=1.4em]
  \item \textbf{სწრაფი მეტრიკების ბარათები (Summary KPIs)}: ცხრილისა და
    გრაფიკის თავზე გამოისახება 4 საკვანძო ინდიკატორი -- ჯამური შემთხვევები,
    საშუალო ტემპერატურა (მინ/მაქს დიაპაზონით), ჯამური/საშუალო ნალექი და
    პიკური აფეთქების პერიოდი შემთხვევების რიცხვით;
  \item \textbf{ინტერაქტიული გრაფიკის ჩამრთველები (Dynamic Series Toggles)}:
    მომხმარებელს შეუძლია ჩამრთველი ღილაკებით (FilterChips) რეალურ დროში
    ჩართოს ან გამორთოს სასურველი მრუდები (შემთხვევები, ტემპერატურა, ნალექი);
  \item \textbf{ორღერძიანი ადაპტური მასშტაბირება}: გრაფიკი დინამიურად
    ირჩევს საბაზისო ღერძს არჩეული მეტრიკების მიხედვით, რაც უზრუნველყოფს
    მკაფიო აღქმადობას მაშინაც კი, როცა მხოლოდ ნალექი ან მხოლოდ ტემპერატურაა
    არჩეული.
\end{itemize}

%=====================================================================
\section{ამჟამად განხორციელებული ფუნქციები}

\begin{itemize}[leftmargin=1.4em]
  \item სტრუქტურირებული ფორმა (დაავადება, რეგიონი, ცვლადები, პერიოდი,
    წყარო), პლუს არასავალდებულო ბუნებრივი ენით მოთხოვნის აღწერა,
    რომელსაც LLM ფორმის წინასწარი შევსების სახით ამუშავებს
    (იხ. \S\ref{sec:nl-parsing});
  \item 39 რეგიონი რეალური, თვიური სიხშირის დენგეს მონაცემებით;
  \item სამი ალტერნატიული სამეცნიერო/ოფიციალური კლიმატის წყარო
    (Open-Meteo, NASA POWER, TMD -- ბოლო მხოლოდ ტაილანდისთვის);
  \item დაავადების მონაცემთა წყაროების ცოცხალი ძებნა WHO GHO-სა და
    HDX-ში, ერთი ჩაშენებული წყაროს ნაცვლად (იხ. \S\ref{sec:source-discovery});
  \item AI-ზე დაფუძნებული მონაცემთა ჰარმონიზაცია, სვეტების მოქნილი ამოცნობა
    და გამჭვირვალე Audit Log ინტერაქტიული მოდალური ფანჯრით (იხ. \S\ref{sec:data-harmonization});
  \item მრავალმრუდიანი ინტერაქტიული გრაფიკი ჩამრთველებით (შემთხვევები,
    ტემპერატურა, ნალექი) და KPI ინდიკატორების ბარათები (იხ. \S\ref{sec:interactive-visualization});
  \item მომხმარებლის მიერ მითითებული URL ან ატვირთული ფაილი, როგორც
    დამატებითი მონაცემთა წყარო;
  \item რეგიონის არჩევა ინტერაქტიული 2D რუკით, სადაც თითოეულ
    მხარდაჭერილ რეგიონს მარკერი აქვს ზუსტ საცნობარო კოორდინატზე;
  \item ყოველი ნაბიჯის ტექსტური ჩვენება მომხმარებელს, და გამჭვირვალობა
    იმის შესახებ, ახსნა ნამდვილად LLM-მა გენერირა თუ fallback ტექსტია;
  \item ცნობილი, მოსალოდნელი backend შეცდომების (მაგ. TMD-ის
    კონფიგურაციის ან custom წყაროს ფორმატის პრობლემა) მკაფიო 502
    პასუხად დაბრუნება ბუნდოვანი 500-ის ნაცვლად;
  \item შედეგის კეშირება \emph{last\_verified} დროის ნიშნულით;
  \item Gemini API-ს გამოძახებების სიხშირის შეზღუდვა (rate limiting);
  \item მიღებული მონაცემების გადმოწერა CSV/JSON ფორმატში და სრული PDF ანგარიში;
  \item ერთი Docker კონტეინერი, რომელშიც Backend ემსახურება
    frontend-საც -- მარტივი განთავსებისთვის.
\end{itemize}

\section{შემდეგი ნაბიჯები}

\begin{itemize}[leftmargin=1.4em]
  \item TMD-ის რეალურ API პასუხზე smoke-ტესტი და, საჭიროებისამებრ,
    \texttt{\_parse\_tmd\_response}-ის ველების დაზუსტება, როგორც კი
    წვდომის მონაცემები დარეგისტრირდება;
  \item მომავალი კლიმატის პროგნოზების (CMIP6/CORDEX) და bias correction-ის
    ინტეგრაცია ემისიის სცენარების მხარდასაჭერად;
  \item გრიდირებული (რეგიონის მთელ ტერიტორიაზე გავრცობილი) მონაცემები
    ერთი წარმომადგენლობითი წერტილის ნაცვლად;
  \item ნამდვილი ფონური დამუშავების რიგი (job queue) გრძელვადიანი
    მოთხოვნებისთვის;
  \item დაავადების მოდელების პირდაპირი მიბმა გამომავალ მონაცემებზე.
\end{itemize}

\end{document}
